\documentclass[11pt]{article}
\usepackage[utf8]{inputenc}
\usepackage[T1]{fontenc}
\usepackage[spanish,es-nodecimaldot]{babel}
\usepackage[a4paper,margin=2.5cm]{geometry}
\usepackage{amsmath,amssymb,booktabs,graphicx,float,hyperref}
\newcommand{\figroot}{../../outputs/toi_ati_future_validation/figures_pdf}
\title{Validacion futura de TOI/ATI}
\author{Proyecto Mapper/TDA de exoplanetas}
\date{}
\begin{document}
\maketitle

\section{Introduccion}
Esta fase no crea otro indice principal. Su objetivo es validar y estabilizar TOI/ATI para convertir un ranking exploratorio en una priorizacion mas auditada.

\section{Motivacion}
ATI original puede priorizar anclas sensibles al radio porque $\Delta_{\mathrm{rel,best}}$ resume el maximo local. La fase presente separa casos estables de casos sensibles y construye una version conservadora del ranking.

\section{Datos usados}
La validacion consume tablas de TOI/ATI, deficit por radio y anatomia previa, sin recalcular Mapper ni introducir un nuevo pipeline de completitud instrumental.

\section{Estabilidad del deficit por radio}
Se comparan tres escalas locales: $r_{\mathrm{kNN}}$, $r_{\mathrm{node\_median}}$ y $r_{\mathrm{node\_q75}}$. Un caso puede ser \textit{stable\_positive\_deficit}, \textit{radius\_sensitive\_deficit}, \textit{unstable\_due\_to\_large\_radius} o \textit{no\_deficit\_or\_overpopulated}.

\section{ATI conservador}
Se construyen variantes interpretativas: ATI\_stable\_simple, ATI\_radius\_penalized y ATI\_conservative. La version conservadora penaliza sensibilidad al radio y radios grandes negativos.

\section{Resultados: casos estables vs sensibles}
\begin{figure}[H]\centering
\includegraphics[width=.82\textwidth]{\figroot/stable_vs_sensitive_deficit.pdf}
\caption{Casos estables vs sensibles al radio.}
\end{figure}

\begin{figure}[H]\centering
\includegraphics[width=.82\textwidth]{\figroot/deficit_profiles_selected_anchors.pdf}
\caption{Perfiles de deficit relativo por radio para anclas seleccionadas.}
\end{figure}

\section{Resultados: cambio de ranking}
\begin{figure}[H]\centering
\includegraphics[width=.82\textwidth]{\figroot/ati_vs_ati_conservative.pdf}
\caption{ATI original vs ATI conservador.}
\end{figure}

\begin{figure}[H]\centering
\includegraphics[width=.82\textwidth]{\figroot/rank_shift_after_stability_penalty.pdf}
\caption{Cambio de ranking tras penalizar sensibilidad al radio.}
\end{figure}

\section{Casos finales para trabajo futuro}
La tabla \texttt{final\_future\_work\_cases.csv} resume cinco casos: region top por TOI, ancla top por ATI original, ancla top por ATI conservador, ancla repetida y ancla con deficit estable.

\begin{figure}[H]\centering
\includegraphics[width=.9\textwidth]{\figroot/final_future_work_cases.pdf}
\caption{Casos finales para trabajo futuro.}
\end{figure}

\section{Prioridad observacional}
La prioridad observacional no implica confirmaciones de objetos ausentes. Es una priorizacion prudente para inspeccion futura bajo referencias topologicas locales.

\begin{figure}[H]\centering
\includegraphics[width=.82\textwidth]{\figroot/observational_priority_ranking.pdf}
\caption{Ranking de prioridad observacional futura.}
\end{figure}

\section{Discusion}
Esta fase fortalece algunos casos estables y vuelve mas fragiles otros. Si un caso como HIP 97166 c baja al penalizar sensibilidad al radio, eso indica que su lectura depende de la escala local. Si HIP 90988 b o HD 4313 b suben por estabilidad, eso sugiere una prioridad mas defendible para inspeccion futura.

\section{Limitaciones}
No hay funcion de completitud instrumental real, $R^3$ simplifica la fisica, el proxy RV no es amplitud real, Mapper solapa membresias y el deficit depende de la referencia local. Los rankings son herramientas de priorizacion, no deteccion.

\section{Trabajo futuro}
\subsection{Incorporar funciones de completitud instrumental}
Modelar $P(\mathrm{detectado}\mid M_p, P, a, estrella, metodo)$ y usar curvas de sensibilidad por metodo.

\subsection{Validar con catalogos futuros}
Correr el pipeline con versiones futuras del catalogo y observar si regiones TOI altas reciben nuevos planetas o mediciones mas finas.

\subsection{Inyeccion-recuperacion sintetica}
Insertar planetas sinteticos en $R^3$, simular seleccion por metodo y evaluar si TOI/ATI recupera regiones incompletas.

\subsection{Propiedades estelares}
Agregar masa, radio, temperatura, magnitud y distancia para mejorar el contexto fisico y la detectabilidad dinamica.

\subsection{Comparacion por metodo de descubrimiento}
Construir rankings separados para RV, transito, imaging y microlensing.

\subsection{Reporte de candidatos observacionales}
Producir fichas para top 5 candidatos y separar casos robustos de casos exploratorios.

\section{Conclusion}
La validacion futura convierte TOI/ATI de un ranking exploratorio en un sistema de priorizacion auditado: no detecta planetas ausentes, pero identifica donde el catalogo parece mas informativo para buscar incompletitud observacional.
\end{document}
